\section{Small-World Routing}
\label{sec:smallworld}

The mechanism by which a request for contract $\contract$ finds peers near $\ell(\contract)$ is small-world greedy routing on the ring. This section describes the topology --- how connections come into existence --- and the routing rule that uses it; the adaptive variant of the rule is the subject of \cref{sec:adaptive}.

\subsection{Background: Why Small-World}

In a regular topology where each peer is connected only to its immediate neighbors on the ring, greedy forwarding visits $\Theta(n)$ peers to traverse a constant fraction of the ring. In a uniformly random topology, paths are short ($O(\log n)$) but greedy forwarding using only local information fails: a peer has no reason to think any particular neighbor is closer to the target. Kleinberg's small-world result \citep{kleinberg2000small} reconciles these by introducing long-range links whose lengths are sampled from a power law: with link probability proportional to $1/d^{r}$ on a $r$-dimensional lattice, greedy routing succeeds in expected $O(\log^2 n)$ hops. Freenet inherits the intuition --- mostly local connections, with a sprinkling of long-range ones, navigated greedily --- but constructs the topology by a different mechanism.

\subsection{Topology by Acceptance Rule}

Each peer chooses its location deterministically at startup as $\ell(p) = H(\mathrm{addr}_p) \bmod 1$. A peer $p$ joining the network sends a \textsf{CONNECT} message to a known \emph{gateway} peer, with the target field set to $\ell(p)$ itself. The gateway, and each subsequent peer that receives the message, applies the following rule:

\begin{quote}
\emph{If I have a neighbor closer to the target than I am, forward the \textsf{CONNECT} to that neighbor; otherwise accept the connection.}
\end{quote}

We call this the \emph{accept-only-at-terminus} rule. It has three consequences worth stating explicitly:

\paragraph{Local clustering.} The peer that finally accepts the connection is, by construction, near the new peer's location on the ring. Joining therefore reliably creates a connection between two peers that are close on the ring. Repeated over many joiners, this fills out the local neighborhood of every peer.

\paragraph{Long-range links from ordinary traffic.} \textsf{CONNECT} is not the only message that traverses the ring. Every \textsf{PUT}, \textsf{GET}, \textsf{UPDATE}, and \textsf{SUBSCRIBE} is forwarded greedily toward a contract location chosen by application logic, and \emph{any} of these requests, when terminating, can leave behind a fresh connection between the originating peer and a peer at an essentially arbitrary ring distance. The application-driven distribution of contract locations is the source of long-range links; we do not need to sample link lengths from a distribution.

\paragraph{Bounded fan-out.} A peer accepts only when it cannot forward closer; once its neighborhood is densely packed, the local terminus typically lies elsewhere. The connection set is further bounded between configurable minimum and maximum (typically 25 and 200 in deployed configurations) by ordinary admission control.

\subsection{Greedy Forwarding}

Routing of an ordinary request from a peer $p$ to a target location $t$ proceeds as follows. Let $\mathrm{N}(p)$ be $p$'s current set of neighbors. If some neighbor $q \in \mathrm{N}(p)$ is strictly closer to $t$ on the ring than $p$ is itself, forward the request to the closest such $q$; otherwise $p$ is a terminus for the request and serves it locally.

Each request carries a \emph{hops-to-live} counter (HTL) that decrements on each forward and bounds the path length; deployed defaults are $\mathrm{HTL}_{\max} = 10$. Above a configurable threshold (HTL $> 7$) the next hop is chosen by random walk rather than greedily; this fallback exists to ensure the request leaves the immediate neighborhood of a peer that happens to be very near $t$ but has not yet cached the contract.

\subsection{Routing Performance}

Small-world topologies built by accept-only-at-terminus exhibit the properties greedy routing requires:

\begin{enumerate}[leftmargin=*, itemsep=2pt]
  \item \emph{High clustering.} Most of a peer's connections are to nearby peers on the ring, because nearby peers are the typical terminuses of join-requests targeted near each other's locations.
  \item \emph{Short paths.} Long-range links exist because peers route requests to contracts at arbitrary locations, and any of those routes can drop a connection at a long-range terminus.
  \item \emph{Robustness.} No single peer is structurally critical: every peer is the terminus of approximately the same number of nearby joiners, and long-range links are diffused across the entire ring.
\end{enumerate}

The expected hop count for greedy routing to a target location is logarithmic in network size in regimes where the long-range link distribution approximates Kleinberg's $1/d$ law. We have not attempted a formal proof of this for the accept-only-at-terminus construction; the empirical hop counts observed in deployment and in simulation are consistent with $O(\log n)$ behavior at the scales we have measured.

\subsection{Bootstrap}

The set of gateway peers is configured out-of-band. A joining peer initiates \textsf{CONNECT} to one or more gateways and continues with the acceptance rule above; the first peer along the route to observe the joiner's external UDP address acts as its NAT discovery proxy. Once the new peer has a small set of neighbors, the bootstrap phase ends and ordinary request traffic drives further topology evolution.
